\section{Liberalism and the Individual}

We now resume our review of one of the newer works of political theory. The point is not to critique it, but rather to
see how its analyses and prescriptions expand or improve on similar projects based solely on Tradition. In \textit{The Fourth
Political Theory}, \textbf{Alexander Dugin} defines the essence of liberalism:

\begin{quotex}
Liberalism as a whole rests on the individual as its most basic component. It is these individuals, collectively but in
isolation from one another, that are taken as the whole. 

\end{quotex}
Of course, we have previously remarked on this idea. In this view, society is a ``heap" rather than an organic whole.
That is why the liberal can sincerely believe that wholesale population replacement will result in no essential change
to the liberal society. The corollary of the liberal assumption is clearly made by Dugin:

\begin{quotex}
Liberalism as an ideology [calls] for the liberation from all forms of collective identity in general. 

\end{quotex}
In particular, this involves the ``denial of \emph{ethnos}", hence liberalism is ultimately an ethnocide. Moreover, other
forms of collective identity are then seen as forms of oppression; this include race, gender sexual relations, and
finally even gender itself. They are seen as artificially constructed categories that are incompatible with the freedom
and full self-expression of the individual.

Dugin takes liberalism as the ``first political theory", at least if we are counting from the period called the
Enlightenment. The second and third theories, intended to oppose liberalism, were communism and Fascism. Theoretically,
these movements tried to supplant the individual as the basic atom of society with either class, race, or state.
Nevertheless, liberalism triumphed over all of them. From the Traditional point of view, the second and third theories
were just as inadequate as the first, since they ignore the fundamental understanding of transcendence. Dugin asserts:

\begin{quotex}
Man is anything but an individual. 

\end{quotex}
The Traditional alternative is that man is also a Person, to use the Medieval idea, also used by Guenon and Evola. Dugin
plans to use the notion of Dasein as developed by \textbf{Martin Heidegger} as the model to oppose the individual of
liberalism. Yet, another existentialist, \textbf{Jean Paul Sartre} brought to light the fundamental presuppositions of
liberalism, both of which are contradicted by Tradition. These are:

\begin{itemize}
\item Existence precedes essence 
\item There is no transcendental ego 
\end{itemize}
In this view, man exists first then chooses in freedom who he is, i.e., his essence. Hence, he can be whatever he
chooses to be. The self is not transcendent to the world, but is somehow embedded in it, and therefore limited by its
``facticity". In liberal social policy, this facticity must be overcome either through social engineering or
technological means. For example, the liberal may believe that anyone can become a ``rocket scientist" and will base
education policy on this belief. The reality of the dearth of rocket scientists can be overcome by ``better schools",
better teachers, better methods, and so on. The one exception who becomes a rocket scientist out of a hundred is
regarded with pride and a vindication of the policies rather than the adventitious event that it is.

By ignoring transcendence, the liberal is subject to the two primal soul or psychic forces: \emph{eros} and
\emph{thumos}. Hence, the idols of modernity are the pop star and athlete, who exemplify the cancerous growth of
\emph{eros} and \emph{thumos} respectively. In contrast, in the Medieval era, the primary types according to
\textbf{Frithjof Schuon} (\textit{Light on Ancient Worlds}), were the Saint, the Hero, and the Sage who exemplifies the
Intellect, the transcendent power that dominates the soul forces.

The Person, as such, is transcendent, yet takes on an essence when manifesting in space-time. In contrast to
existentialism, the Person is transcendent to existence, yet has essential characteristics. In this view, Man is
related in two directions: vertically, to the Absolute via transcendence, horizontally, to others through his essential
characteristics of gender, family, and ethnos. Yet these serve to fulfill the Person, not to oppress him. Without these
larger supports, the individual as such can achieve little. Dugin writes:

\begin{quotex}
The individual is granted freedom because the uses to which he can put it are extremely limited—it
will remain contained within the tiny scope of his individuality and that over which he has direct control. This is the
flip side of liberalism: at its core, it is totalitarian and intolerant of differences, and most especially opposed to
the realization of a great will. It is only prepared to tolerate small people; it protects not so much the rights of
man, but, rather, the rights of a small man. This small man can be allowed to do anything, but in spite of all his
desire, he will be unable to do anything. 

\end{quotex}
This is true. For many ``individuals", the most creative act of their life may be no more than choosing a tattoo or
learning a new sex position. They may have little awareness of themselves as transcendent and fail to develop their
Intellect. The Traditional society is criticized for being patriarchal and hierarchal; that is because no spiritual or
political authority is recognized as natural, but rather authority depends solely on the consent of the individual.
Nevertheless, the intent is not to oppress the masses, but rather to help them develop in the correct direction. Those
with a weak Intellect or Will, can be guided and directed by the stronger members.

Some may object to this description on the grounds that no ``liberal" truly acts the way described. We are merely
pointing out the axioms and logic of liberalism. If no one can consistently and coherently think and act that way, then
that is the limitation of the ideology itself. Of course, few will live, even in the modern age, solely by the
unchecked forces of \emph{eros} and \emph{thumos}, and they will bring their Intellect into their way of life, even if
not fully aware of it.



\flrightit{Posted on 2012-10-11 by Cologero }

\begin{center}* * *\end{center}

\begin{footnotesize}\begin{sffamily}



\texttt{Mihai on 2012-10-11 at 03:46 said: }

Liberalism is tolerance for all and everything EXCEPT for that which opposes ``anything goes" and its insane
consequences. 

``at its core, it is totalitarian and intolerant of differences"

True, in liberalism the individual is only granted an external and chaotic type of freedom. Internally, he is to be
subjected to an exterior and alien role, a role which society has come to name ``everyday life", or ``real life". It
never even occurs to anyone to question the meaning of ``real" in this context, because the power of liberalism is not
based on principles and logic, but merely on jargon and slogans. 

So, while ``anything goes", a truly meaningful existence and lifestyle has become a luxury, which the liberal
establishment is going to allow, if at all, only when the individual has fulfilled his obligation to ``normal, everyday
life".


\hfill

\texttt{Dominion on 2012-10-11 at 04:00 said: }

A mistake I think Dugin makes in his fourth theory is that he rejects almost out of hand the reasons why liberalism
arose. Liberalism makes a valid point, though from false principle, in that legitimate authority must be for the good
of those which it has power over. This is easy to say, but in practice extremely difficult to achieve, particularly on
the social and political plane. ``For the good of the people" has been the rallying cry of many a despot and tyrant. It
seems that those who can truly claim legitimate Traditional authority will rarely explicitly state that they have it,
much less demand relentlessly that they be recognized as such. Hierarchy is a value which needs defending; the problem
is that many who defend it have a tendency to use valid principles for their own ends.

Dugin's use of being-in-the-world as a basis for an ideology is far more ingenious than some seem
to realize. While it is true, as some have stated, that the people will likely not be `storming
the Bastille' anytime soon over `Dasein', However, the
fourth political theory must be one which creates a slow shift, an enduring force in this waning of the current cycle,
a theory which can be applied to people living their lives, formed by ethny and gender and culture and transcendence;
not simply another way for the angst-ridden to try and realize their dream of a sexy revolution, flags and thunder
behind them. Purposefully applying the principles of Tradition to shape ones being-in-the-world as a person, a people,
and a civilization seems a prime way to allow Tradition to manifest in countless forms. 

A question to you, Cologero: do you look at ethny and gender, etc, as being forms which allow a person to realize their
transcendent Personhood, or are these aspects of the transcendent Person? Also, is this Person identical to the ``true
Self", or Atman?


\hfill

\texttt{william zeitler on 2012-10-11 at 10:05 said: }

I've been following this blog for maybe a year, and still have only a vague idea of what you all
mean by `Tradition'. Is there a past post or article somewhere that clearly
defines/explains what YOU mean by this term? Thanx!


\hfill

\texttt{Cologero on 2012-10-21 at 14:26 said: }

Here is an essay, \textit{Big Bird, Liberalism, and Perversion}\footnote{\url{http://www.americanthinker.com/2012/10/big_bird_liberalism_and_perversion.html}}, that understands that the aim of liberalism is to obliterate all
distinctions. Unfortunately, like most modernists, the author takes a medical approach and considers liberals as
suffering from a childhood developmental disorder. If that were true, the antidote to liberalism would be some
technique of psychotherapy or perhaps even a new psychoactive drug.


\hfill

\texttt{Logres on 2012-10-21 at 22:34 said: }

Liberalism is always willing to annihilate actual inner diversity (usually built on history, race, culture, religion or
a combination of such) in favor of a faux outer diversity; the price for this seems to be a total inner conformity
and/or emptiness, itself a result of the ruins of that upon which Liberalism waded to power. I often wonder why it
never occurs to them that there are trade-offs in reality, although I can see why their elite master class pursues it.
The one thing they are absolutely opposed to is external monopoly (even if it protects an inner diversity or itself
forms part of a larger global patchwork of diverse forms of government), yet this is precisely where they end up
anyway: in the end, they will lose the material diversity they sold their soul for, after all. In any case, they have
yet to develop a coherent, positive basis for discriminating between good and evil so that both are not equally
``tolerated". Hence, chaos and lies and confusion abound. 

Or so it seems to me lately…and I think about it, over and over again, trying to see some good way out for the Liberal
regime…


\end{sffamily}\end{footnotesize}
